\chapter{Discusión y análisis de resultados}
\label{cap:discusion}

% =====================================================================================
% CAPITULO NUEVO (ago-2026). Corresponde al capitulo 6 de la estructura del tipo 3:
% «Discusion y analisis de resultados».
%
% ⚠️ LA DIVISION DE TRABAJO CON EL CAPITULO 10 NO ES NEGOCIABLE. La norma dice del
%    capitulo de resultados: «Es una exposicion objetiva, SIN valorar los resultados ni
%    justificarlos». Toda la interpretacion vive AQUI.
%
% ⚠️ La rubrica pesa este capitulo mas que ningun otro: «Desarrollo especifico de la
%    contribucion» vale el 20 %, y nombra «Discusion correcta de los resultados» en las
%    bandas de notable y sobresaliente.
%
% CAPITULO BLOQUEADO hasta que existan los resultados del capitulo 10.
% =====================================================================================

\section{¿Aporta algo la estructura del circuito?}
\label{sec:aporta_discusion}
% LA pregunta del TFM, planteada en el capitulo 9 y respondida aqui.
% Interpretar la diferencia entre el GEM y los dos tabulares: si existe, a que se debe;
% si no existe, que dice eso sobre el problema.

\section{Ventajas y desventajas de cada solución}
\label{sec:ventajas}
% La norma del tipo 3 lo pide con estas palabras: «el analisis de las ventajas y
% desventajas de las distintas soluciones evaluadas». Una subseccion por modelo, y no
% solo en terminos de error: coste de entrenamiento, interpretabilidad, y comportamiento
% fuera de distribucion.
\subsection{Ridge}
\subsection{Random Forest}
\subsection{El \textit{Graph Transformer}}

\section{Qué transfiere y qué no, eje por eje}
\label{sec:transferencia}
% Enlazar con el limite honesto del capitulo 8: en validacion funcionaba y en test solo
% transferia un observable. ¿Cambia eso con el modelo de grafo?
\subsection{Generalización a tipos de circuito nunca vistos}
\subsection{Extrapolación en tamaño}
\subsection{Robustez frente a la deriva del hardware}

\section{Datos anómalos y casos de fallo}
\label{sec:anomalos}
% La norma menciona explicitamente «presentar posibles explicaciones para los datos
% anomalos». Candidatos ya conocidos: las muestras sin señal que perder, las colas
% pesadas que dominan el MAE, y el label shift de std_Z.

\section{Contraste con el estado del arte}
\label{sec:contraste}
% Volver al capitulo 4 y cerrar el circulo: los cuatro diferenciadores anunciados alli,
% ¿se sostienen con los numeros en la mano?

\section{Amenazas a la validez}
\label{sec:validez}
% Distinto de las limitaciones del capitulo 12: aqui van las que afectan a la LECTURA de
% estos resultados concretos.
%   - el modelo de ruido no incluye crosstalk;
%   - el eje de tamaño llega hasta 15 qubits, no mas;
%   - un solo backend simulado, no varios.
